\begin{center}
\begin{tikzpicture}[x=.9cm,y=.9cm,font=\small]
  \fill[paper] (0,0) rectangle (10.8,5.8);
  \draw[gray!55,line width=1pt] (3.6,.5) .. controls (4.3,1.7) and (3.7,3.7) .. (4.5,5.4);
  \node[rotate=76,text=gray!70] at (3.95,4.1) {崤函山地};
  \draw[blue!55,line width=1pt] (.4,4.6) .. controls (3.2,4.25) and (6.7,4.95) .. (10.2,4.45);
  \node[text=blue!70!black,above] at (6.7,4.65) {黄河方向};
  \draw[blue!45,line width=.8pt] (4.6,.5) .. controls (5.0,1.4) and (5.6,2.35) .. (5.5,3.65);
  \node[text=blue!70!black,right] at (5.5,1.35) {伊水、洛水方向};

  \fill[dynasty!24] (.8,1.2) -- (3.25,1.0) -- (3.45,4.55) -- (1.25,4.9) -- cycle;
  \node[align=center] at (2.05,2.95) {秦\\关中};
  \fill[yellow!24] (4.15,2.75) -- (5.15,2.6) -- (5.35,3.75) -- (4.35,3.95) -- cycle;
  \node at (4.75,3.28) {韩};
  \fill[green!20] (5.85,3.0) -- (8.25,2.9) -- (8.45,4.55) -- (6.05,4.7) -- cycle;
  \node at (7.15,3.78) {魏};

  \fill[timeline] (3.85,2.95) circle (.12);
  \node[below,align=center] at (3.85,2.88) {函谷\\方向};
  \fill[timeline] (5.25,2.05) circle (.14);
  \node[below,align=center] at (5.25,1.96) {伊阙\\前293年};
  \fill[source] (5.62,3.72) circle (.1);
  \node[above,align=center] at (5.62,3.74) {洛阳\\方向};

  \draw[-{Stealth[length=2mm]},ultra thick,color=dynasty] (2.9,2.65) -- (5.05,2.1);
  \draw[-{Stealth[length=2mm]},thick,color=yellow!60!black] (4.68,2.75) -- (5.12,2.2);
  \draw[-{Stealth[length=2mm]},thick,color=green!60!black] (6.22,3.05) -- (5.4,2.2);
  \node[align=center,text=dynasty] at (3.55,1.92) {秦军东出\\与补给方向};
  \node[align=center,text=source] at (7.85,2.25) {韩、魏联军\\向峡口集结};
\end{tikzpicture}

\smallskip
\small\textit{图\quad 伊阙—函谷走廊示意：据《史记·秦本纪》《白起王翦列传》及《资治通鉴》呈现前293年伊阙之战所处的关中东出口、伊洛峡口与韩魏防线之间的关系。山地、河流和地点用于说明通道，不复原军队日次部署或补给线；各政治体边界均为约略示意。}
\end{center}
